%PATH=/usr/local/texlive/2019/bin/x86_64-linux:$PATH

\documentclass[letterpaper,twocolumn,openany, twoside, 11pt, usenames]{cocbook}

\usepackage[utf8]{inputenc}
\usepackage[english, russian]{babel}
\usepackage{noto}
\usepackage{fontspec}
\usepackage{titlesec}
\usepackage{lipsum}
\usepackage{hyphenat}
\usepackage{graphicx}
\usepackage{enumitem}
\usepackage[table]{xcolor}
\usepackage{tabularx}
\usepackage{multicol}
\usepackage{cuted}

\newfontfamily\headingfont[]{Oswald-Regular.ttf}
\newfontfamily\smallheadingfont[]{NotoSerif}
\setmainfont[]{NotoSerif}
\titleformat{\section}
  {\headingfont\fontsize{20}{1em}\bfseries\color{RedCthulhu}}
  {\MakeUppercase{\chaptertitlename} \thesection:}
  {1em}
  {\MakeUppercase}

\titleformat{\subsection}[hang]
  {\smallheadingfont\fontsize{14}{1em}\bfseries\color{RedCthulhu}}
  {\MakeUppercase{\chaptertitlename}\ \thesubsection:}
  {1em}
  {\MakeUppercase}
\titleformat*{\subsubsection}{\normalfont\fontsize{12}{1em}\bfseries\color{RedCthulhu}}

\renewenvironment{quote}{%
   \list{}{%
     \leftmargin0.0cm   % this is the adjusting screw
     \rightmargin\leftmargin
   }
   \item\relax
}
{\endlist}

\let\oldquote\quote
\let\endoldquote\endquote
\renewenvironment{quote}[2][]
  {\if\relax\detokenize{#1}\relax
     \def\quoteauthor{~#2}%
   \fi
   \oldquote}
  {\par\nobreak\smallskip\hfill(\quoteauthor)%
   \endoldquote\addvspace{\bigskipamount}}

\hyphenchar\font=-1
\sloppy

\newcolumntype{Y}{>{\centering\arraybackslash}X}

\textwidth=6.70in
\linespread{1.5} 

\begin{document}
\scriptsize

\section*{\nohyphens{Что такое зов Ктулху?}}
\begin{quote}{--- Г.Ф. Лавкрафт}
{\it Самая старая и самая сильная эмоция человечества --- это страх. А старейшим и сильнейшим видом страха является страх перед неизвестностью.}
\end{quote}

\begin{quote}{Г.Ф. Лавкрафт, --- {\it Безымянный город}}
{\it То не мертво, что вечность охраняет, Смерть вместе с вечностью порою умирает.}
\end{quote}

{\it Зов Ктулху} --- игра о тайнах, загадках и ужасе. Игроки берут на себя роль отважных исследователей и путешествуют в странные и опасные места, раскрывая зловещие заговоры и защищая мир от ужасов ночи. На своем пути они встретят существ, обитающих за пределами пространства и времени, само существование которых угрожает рассудку, омерзительных чудовищ и безумных культистов. Погружаясь в тексты старинных фолиантов, им предстоит раскрывать секреты, которые не должен знать ни один человек. Они встретят множество преград на своем пути, но их героические действия вполне могут определить судьбу мира.

Созданная Сэнди Питерсеном и впервые опубликованная в 1981 году настольная ролевая игра {\it Зов Ктулху} на протяжении более 35 лет определяла жанр и часто называется одной из лучших ролевых игр. Тем, кто достаточно смел духом, чтобы раскрыть ее тайны, эта игра готовит награды за пределами любых ожиданий!

В игре каждый игрок берет на себя роль персонажа, а один из игроков - роль рефери --- Хранителя Тайного Знания (``Хранителя''), который служит в игре модератором и обрисовывает игрокам сюжет и сеттинг, в рамках которых проходит их приключение. Используя игровые кубики и несложные правила, вы определяете успех или неудачу действий персонажей, в то время как сюжет заставляет их попадать в сложные и опасные ситуации.

Этот буклет позволит вам создать персонажа для настольной ролевой игры {\it Зов Ктулху}, а также обучит основам игровых правил, достаточных для того, чтобы вы с вашими друзьями могли начать играть и насладиться стартовыми приключениями. Полноценные книги правил, а также приключения можно купить на сайте создателей игры --- компании Chaosium.

\subsection*{\nohyphens{Роли игроков}}

В каждой игре игроки берут на себя одну из двух ролей: исследователя либо Хранителя. Большинство выбирает роль исследователя, и исследование --- это основное, чем они будут заниматься: решать загадки или искать выход из сложных ситуаций. Сюжеты будут бросать персонажам вызов: они могут оказаться ранеными, испытать неподвластные рассудку видения или даже быть проглоченными чудовищем! По мере того, как игра будет развиваться, исследователи смогут узнать о странной магии и ужасных, чуждых созданиях, получить тайные знания из магических гримуаров, полных зловещих тайн, и улучшить свои навыки, становясь более опытными и готовыми к приключениям.

Один из игрокок берет на себя роль Хранителя. Он выбирает сценарий для игры или может придумать собственный оригинальный сюжет. В игре Хранитель подготавливает сцену, описывает происходящее и изображает людей, которые попадаются исследователям на пути (они обычно называются NPC, от английского ``Non-Player Character'' --- персонаж, роль которого исполняет Хранитель). Хранитель также помогает определять результат действий и служит арбитром в части игровых правил. Так как Хранитель должен тратить дополнительное время и силы на подготовку к игре, роль Хранителя от сценария к сценарию переходит от игрока к игроку. О Хранителе можно думать как о режиссере фильма, актеры которого (игроки) не знают, что будет дальше по сюжету.

Сам процесс игры представляет собой развивающееся взаимодействие между игроками --- в образе их персонажей, разгадывающих тайну, --- и Хранителем, который описывает мир этой тайны.

В этой игре нет никакого игрового поля. Процесс представляет собой простой разговор: Хранитель моделирует некую ситуацию, а затем игроки от лица своих персонажей рассказывают о том, что они хотят сделать.

Используя систему правил, чтобы все было честно и последовательно, Хранитель сообщает им, могут ли они совершить то, что хотят, и если да, что им нужно сделать, чтобы добиться желаемого эффекта. Обычно это подразумевает бросок кубиков, который позволяет определить результат событий, а также гарантирует непредвзятость происходящего и добавляет накала страстей. Результат броска может привести к непредвиденному сюрпризу, жестокому поражению или, напротив, позволить персонажу избежать смерти в последнюю секунду! После того, как результат броска определен, Хранитель описывает происходящее и спрашивает у игроков, как реагируют на это их персонажи. Так и проходит игра.

Цель ролевой игры - повеселиться. Даже когда сердца бьются, а со лба капает пот, людям нравиться пугаться, если только ужасы не настоящие. Для некоторых расслабление после пережитого страха --- само по себе желанный результат. Другим просто нравится пребывать в этом состоянии. {\it Зов Ктулху} сможет порадовать и тех, и других!

\subsection*{\nohyphens{Игровые концепции}}

{\it Зов Ктулху} основан на историях писателя Говарда Филлипса Лавкрафта (20 августа 1890 - 15 марта 1937), чьи истории имели в своей основе оригинальную мрачную философию, согласно которой, мироздание безжалостно и им правят чуждые  человеку и обладающие богоподобной мощью создания. Эта концепция была достаточно новаторской, чтобы увлечь поколения писателей, развивших и расширивших лавкрафтовскую вселенную, которая получила название ``Мифология Ктулху''.

В течение жизни Лавкрафт был практически неизвестен и публиковался исключительно в дешевых бульварных журналах. Он умер в нищете, но в наше время считается одним из наиболее влиятельных авторов, творивших в жанре ужасов. Несмотря на его устаревшие и во многом расистские взгляды, мы можем вдохновляться творчеством Лавкрафта и использовать его как основу для ролевых игр. Его произведения написаны тяжелым, архаичным языком, как будто сошедшим со страниц тех древних книг, которые часто упоминаются в его рассказах. Однако мы можем заимствовать из его творчества не только ужасы, но и юмористическую сторону: Лавкрафт зачастую подтрунивал сам над собой и над своими коллегами по перу, пародируя их стиль. Современные авторы книг и игр используют творчество Лавкрафта и других его современников, вплетая в повествование как современные, так и исторические проблемы.

Если вы хотите узнать больше о Лавкрафте и Мифологии Ктулху, вашим лучшим помощником станет интернет. Начать можно с Википедии.

\subsubsection*{\nohyphens{Мифология Ктулху}}

Авторство термина ``Мифология Ктулху'', как правило, присваивают Огасту Дерлету --- писателю и раннему поклоннику творчества Лавкрафта. Он стал основателем издательского дома ``Arkham House'', который поставил своей целью сделать широко доступными книги Лавкрафта. Сегодня под этим термином подразумевается вымышленная космология, которая включает богов и чудовищ, тайные знания и то, что называют ``лафкрафтовскими ужасами''. В этой космологии человечеству отведена невзрачная роль быть на задворках вселенной.

Лавкрафт писал: ``Все мои рассказы основаны на предположении, что обычные человеческие законы, интересы и эмоции не имеют никакого значения в контексте мироздания''. Он также полагал, что основополагающие истины вселенной настолько чужды и ужасающи для человеческого восприятия, что малейшее их воздействие на человека может привести к потере разума. И хотя человечество хочет одновременно и покоя, и истины, в космологии Лавкрафта можно обладать только чем-то одним. Человеческий разум --- несовершенное вместилище и не может одновременно содержать космические истины и неповрежденный рассудок.

В Мифологии Ктулху, те люди, что жаждут власти могут избрать путь во тьму и полностью лишиться рассудка в обмен на овладение умением повелевать тайнами времени и пространства. После заключения сделки с дьяволом, эти беспощадные чародеи могут призвать разрушения, которых этот мир еще не видел в обмен на еще большую власть и знания. Чуждые существа Мифологии зачастую безразличны к человечеству, которое часто воспринимает их как божеств и основывает культы, почитающие их. Подобные культисты - одни из основных антагонистов в этой игре.

\subsubsection*{\nohyphens{Тайна и Исследование}}

Приключения в {\it Зове Ктулху} обычно разворачиваются вокруг тайны, а к персонажам-исследователям обращаются, чтобы добраться до правды. Зачастую загадка была порождена преступлениями безумного идолопоклонника или культа в подчинении божеств Мифологии. Задача исследователей - найти улики, каждая из их которых укажет на дополнительные ответы на стоящие перед персонажами вопросы или, вероятно, поставит перед ними новые. По мере того, как исследователи находят все большее количество улик, их понимание ситуации улучшается, пока они не найдут разгадку, которая в свою очередь может привести к изначальному виновнику и финальному противостоянию. Задача игроков, которые отыгрывают этих исследователей, - раскрыть тайну и определиться с тем, как они привлекут виновных к ответственности или расправятся с ними. Иногда тайна может быть результатом странного колдовства, действий омерзительного монстра или другого необычного события, каждое из которых может оказать непредвиденный эффект на исследователей.

Само собой, все приключения разные и не все начинаются с тайны. Иногда сценарий просто начинается сценой, в которой персонажи находятся в центре опасной ситуации из которой им надо выйти. Каждое приключение, как правило, - небольшая история и эти истории отличаются друг от друга. Истории так же могут быть связаны друг с другом, и формировать более масштабный сюжет, который называют кампанией.

\section*{\nohyphens{Создание исследователя}}

Для того, чтобы играть в {\it Зов Ктулху}, вам необходимо создать персонажа - исследователя. Создание персонажа можно представить в виде формальных шагов, они описаны ниже. Игроки записывают информацию о своих персонажах на специальных листах персонажей. Эти листы содержат всю информацию, необходимую для игры. Их можно скачать здесь: https://www.chaosium.com/cthulhu-character-sheets.
Представленный ниже процесс создания персонажа немного упрощен, чтобы дать вам возможность начать играть как можно быстрее. Более подробный процесс описан в {\it Книге Хранителя} и {\it Книге Исследователя}.

\begin{figure}
  \includegraphics[width=\linewidth]{img/cthulhu.jpg}
  \caption{A boat.}
  \label{fig:boat1}
\end{figure}

\begin{figure}
\setthemecolor[CoCPaperBox]
\end{figure}

\subsection*{Шаг первый: характеристики исследователя}
Итак, исследователь в Зове Ктулху имеет восемь характеристик:

\begin{itemize}[leftmargin=4mm]
  \item Сила (Strength, STR, СИЛ) служит мерилом чистой физической мощи вашего исследователя
  \item Телосложение (Constitution, CON, ТЛС) определяет стойкость персонажа и его здоровье
  \item Сила Воли (Power, POW, ВОЛ) - сочетание силы духа, воли и ментальной стабильности
  \item Ловкость (Dexterity, DEX, ЛВК) - мерило физической ловкости и скорости
  \item Внешность (Appearance, APP) отвечает за физическую привлекательность вашего персонажа
  \item Размер (Size, SIZ, РЗМ) отражает сочетание высоты и веса исследователя
  \item Интеллект (Intelligence, INT, ИНТ) - приблизительная прикидка хитроумия вашего персонажа и его способности делать логические и интуитивные выводы
  \item Образование (Education, EDU, ОБР) - аккумулированный жизненный опыт персонажа, будь то результат образования или жизни на улице.
\end{itemize}

Распределите следующие значения как вам угодно среди восьми характеристик: 40, 50, 50, 50, 60, 60, 70, 80. Каждое из этих значение - это процент, т.е. если вы определяете силу персонажа в 70, это означает, что у них ``СИЛ 70\%`` - показатель выше среднего, т.е. персонаж довольно силен. Запишите эти значения в большой ячейке рядом с каждой из характеристик. Эти значения называются ``Обычными''

\subsubsection*{\nohyphens{Половинные и Пятые Значения Характеристик}}

Взглянув на лист исследователя, вы увидите три ячейки рядом с каждой характеристикой: одну большую ячейку (в которой вы только что записали числа для этой характеристики) и еще две маленькие ячейки сбоку. Верхняя маленькая ячейка для ``Половинного'' значения, та что ниже - для ``Пятого''.

\begin{itemize}[leftmargin=4mm]
  \item Для того, чтобы использовать Половинное значение характеристики, просто разделите пополам значение этой характеристики, округляя в меньшую сторону до ближайшего целого значения, если необходимо. Например, если вы выбрали значение СИЛ 70, то Половинное значение (используемое для Сложных бросков) будет равно 35.
  \item Для получения Пятого значения, разделите Обычное значение на 5, также округляя в меньшую сторону. Например, для СИЛ 70 Пятое значение (используемое для Экстремально сложных бросков) будет равно 14.
\end{itemize}

Запишите Половинные и Пятые значения для каждой характеристики в лист исследователя. На стр 23 есть таблица со значениями, позволяющими быстро сориентироваться. Загружаемые PDF листов делают автоматически калькулируют требуемые значения.

\subsection*{Шаг второй: вторичные атрибуты}

Также есть определенное число атрибутов, определяемых после того, как вы определились с характеристиками вашего исследователя. Это бонус к урону (Damage Bonus), комплекция (Build), очки здоровья (Hit Points), скорость передвижения (Move Rate), рассудок (Sanity), очки магии (Magic Points). Кроме этого, вам нужно будет определить показатель удачи.
\smallbreak
\noindent \textbf{Бонус к урону (БУ) и Комплекция:} бонус к урону определяет то, как много дополнительного урона вам исследователь наносит, когда он совершает успешную атаку в ближнем или рукопашном бою. Комплекция - это результат сочетания размера и силы, который обычно используется в ``рукопашных маневрах'' (см. стр. 19). Сложите ваши Обычные значения СИЛ и РЗМ и обратитесь к таблице на стр. 8. На листе исследователя есть ячейки для итогового значения Комплекции и бонуса к урону.
\smallbreak
\noindent \textbf{Очки здоровья (ОЗ):} они являются результатом суммы РЗМ и ТЛС, разделенной на 10 и округленной вниз до ближайшего целого числа. По мере того, как ваш исследователь получает урон в бою или от каких-либо других воздействий, это значение будет уменьшаться.
\smallbreak
\noindent \example{РЗМ 50 и ТЛС 50 в сумме дают 100, разделенное на 10, это значение равно 10 очкам здоровья. Исследователь может получить 10 очков урона, после чего он потеряет сознание и, вполне вероятно, умрет.}
\smallbreak
\noindent \textbf{Скорость перемещения (ПРМ):} все персонажи-люди перемещаются со Скоростью, равной 8.
\smallbreak
\noindent \textbf{Очки рассудка (РАС):} изначально равно показателю ВОЛ персонажа. Вы увидите раздел с очками рассудка на листе персонажа: вам нужно обвести кружком нужное значение. Это значение используется в процентном (1D100) броске, который используется чтобы понять насколько ваш персонаж может сохранять хладнокровие перед лицом абсолютного кошмара. По мере того, как вы сталкиваетесь с ужасами Мифологии Ктулху, ваш показатель рассудка будет изменяться.
\smallbreak
\noindent \example{ВОЛ 40 дает показатель рассудка равный 40. Когда вы делаете бросок на Рассудок, вам нужно выкинуть меньше или равно 40 на 1D100 чтобы бросок считался успешным. Значение 41 или выше означает, что бросок на рассудок провалился.}
\smallbreak
\begin{cocpaperbox}{}{}
  \raisebox{-10mm}{\includegraphics[width=\linewidth]{img/top.png}}
  \cocpaperboxtitle{Таблица бонуса к урону и Комплекции}
  \noindent \example{Рома выбирает СИЛ равную 60 и РЗМ равный 70, в сумме дающие 130. Когда он совершает успешную физическую атаку, он наносит дополнительные 1D4 очка урона (Бонус к Урону). Его комплекция равна 1.}
  \bigbreak
  \rowcolors{2}{}{TableRow}
  \begin{tabularx}{74mm}{@{}YYY@{}}
  \rowcolor{TableHeading}\color{CoCPaperBox}\bfseries СИЛ + РЗМ & \color{CoCPaperBox}\bfseries Бонус к Урону & \color{CoCPaperBox}\bfseries Комплекция \\
  2-64 & -2 & -2 \\
  65-84 & -1 & -1 \\
  85-124 & Нет & 0 \\
  125-164 & +1D4 & 1 \\
  165-204 & +1D6 & 2 \\
  \end{tabularx}
  \includegraphics[width=\linewidth]{img/bottom.png}
\end{cocpaperbox}
\noindent \textbf{Очки магии (ОМ):} равны 1/5 ВОЛ. Обведите в кружок это значение на листе исследователя. Очки магии используются в колдовстве, а также для подпитки магических устройств и волшебных эффектов. Потраченные очки восстанавливаются сами со скоростью 1 очко в час. Если у персонажа не осталось очков магии, все последующие очки вычитаются из его очков здоровья, любая подобная потеря отражается физическим уроном в том виде, в котором реши Хранитель.
\smallbreak
\noindent \example{ВОЛ 40 предоставляет 8 очков магии. При колдовстве заклятия необходимо потратить 2 очка. Количество имеющихся очков у исследователя временно падает до 6.}
\smallbreak
\noindent \textbf{Показатель удачи (Удача):} определите его, кинув 3D6 и умножьте результат на 5. Обведите итоговое значение на листе исследователя. Бросок на показатель Удачи обычно используется, чтобы определить, насколько расположены в вашу или не вашу пользу внешние обстоятельства. Как и с Рассудком, вы должны выкинуть меньше или равно показателю Удачи, чтобы бросок считался успешным. На стр. 15 приведены подробности о броске на Удачу.
\smallbreak
\noindent \example{Персонаж Романа бежит от толпы зомби и запрыгивает в ближайший автомобиль. Хранитель просит Рому сделать бросок Удачи, чтобы определить, находятся ли ключи в зажигании (т.к. это абсолютно случайная вероятность). Роман делает процентный бросок, кидая 1D100 и выкидывает 28, что меньше его показателя Удачи --- он поворачивает ключ и двигатель заводится!}

\subsection*{Шаг третий: род деятельности и навыки}

К этому моменту у вас уже должно быть представление, чем вам исследователь может зарабатывать себе на жизнь. Помните, что ``исследователь'' не обязательно означает палеонтолога или журналиста. Выбор вида занятости определит навыки, доступные вашему персонажу. Для начала нужно выбрать род деятельности. Все, кого вам было бы интересно отыгрывать, подойдет, но итоговый выбор нужно согласовать с Хранителем. Наиболее популярными видами деятельности в {\it Зове Ктулху} являются Профессор, Журналист, Оккультист и Археолог, но в целом ваш выбор ограничен только фантазией.

Либо выберите вашу профессию из приведенного списка и используйте ее навыки, либо, используя общий список навыков создайте свою собственную. Для этого выберите восемь навыков, которые наиболее подходят для вашей профессии. Эти навыки будут ``профессиональными навыками''.
\smallbreak
\noindent \textbf{Примечание:} краткое описание разных навыков может быть найдено на стр. 10-12.
\smallbreak

Теперь вы можете распределить очки навыков на листе персонажа. Ни один персонаж не может добавлять очки к навыку Знания Мифологии Ктулху, т.к. подразумевается, что любой персонаж ничего не знает об ужасах мироздания.

Распределите следующие значения по восьми профессиональным навыкам, а также по навыку Благосостояния: на один из навыков 70\%, на два 60\%, на три 50\% и еще на три 40\% (установите значения выбранных навыков равные именно этим значениям, не обращайте внимания на базовые значения навыков, указанные на листе персонажа).

\smallbreak

\noindent \example{Полина хочет играть за Журналистку и распределяет следующие навыки: Искусство/Ремесло (Фотография) 50\%, История 40\%, Использование Библиотек 50\%, Родной Язык (Английский) 60\%, Психология 40\%. Для общения она выбирает Убеждение 70\% (очень убедительная!). Затем она смотрит на лист персонажа и видит еще два навыка, которые могут оказаться полезными журналисту: Наблюдательность 50\% и Скрытность 60\%. У нее осталось одно значение - 40\%, которое она решает назначить своему Благосостоянию. Полина записывает эти значения в большую ячейку рядом с каждым навыком.}

\smallbreak

После распределения очков по профессиональным навыкам, выберите ``личные навыки''. Это те навыки, которыми ваш персонаж овладел вне своей основной рабочей деятельности. Выберите четыре не профессиональных навыка и улучшите их на 20\% (прибавьте 20\% к значению, указанному для выбранного навыка на листе персонажа).

Мы рекомендуем записывать значения ваших навыков в том же формате, что и Характеристики -  Обычное / Половинное / Пятое - они понадобятся вам во время игры (на стр 23 приведена Памятка со значениями). Само собой, если вам удобнее, можете высчитать эти значения, когда они вам понадобятся во время игры.

\smallbreak

\noindent \example{Рома выбирает в качестве рода деятельности профессию Солдата. Так как этой профессии нет в этой книге, он сам решает, что ему подойдут следующие навыки: Скалолазание, Уклонение, Схватка (Рукопашный Бой), Огнестрельное Оружие (Винтовка/Дробовик), Первая Помощь, Язык (Иной), Скрытность и Выживание. Рома распределяет свои очки следующим образом: Скалолазание 60\%, Благосостояние 40\%, Уклонение 60\%, Схватка (Рукопашный Бой) 70\%, Огнестрельное Оружие (Винтовка/Дробовик) 50\%, Первая Помощь 40\%, Язык (Иной) 50\% (он выбирает Испанский в качестве второго языка), Скрытность 50\%, Выживание 40\%.

Затем Роман выбирает четыре личных навыка (прибавляя 20\% к их стартовому значению на листе персонажа): Вождение Автомобиля 40\%, Механика 30\%, Наблюдательность 45\% и Прыжок 40\%. Каждое значение навыка он записывает в ячейку рядом с соответствующим названием и заодно вносит значения для Обычных/Половинных/Пятых, например ``Наблюдательность 45 (22/9).''
}

\begin{cocpaperbox}{}{}
  \raisebox{-10mm}{\includegraphics[width=\linewidth]{img/top.png}}
  \cocpaperboxtitle{Примеры профессий исследователей}
  \textbf{АНТИКВАР}---Оценка, Искусство/Ремесло (Любое), История, Использование Библиотек, Язык (Иной), один навык для общения (Обаяние, Увещевание, Запугивание или Убеждение), Наблюдательность и один любой навык.
  \bigbreak
  \textbf{ПИСАТЕЛЬ}---Искусство (Литература), История, Использование Библиотек, Знание Природы или Оккультные Знания, Язык (Иной), Родной Язык, Психология и один любой навык.
  \bigbreak
  \textbf{ВОЛЬНЫЙ ХУДОЖНИК}---Искусство/Ремесло (Любое), Огнестрельное Оружие,  Язык (Иной), Верховая Езда, один навык для общения (Обаяние, Увещевание, Запугивание или Убеждение), любые три других навыка.
  \bigbreak
  \textbf{ДОКТОР МЕДИЦИНЫ}---Первая Помощь, Язык (Иной: Латынь), Медицина, Психология, Естественные Науки (Биология), Естественные Науки (Фармакология), любые два других навыка как научные или личные специализации (например, психиатр может выбрать Психоанализ).
  \bigbreak
  \textbf{ЖУРНАЛИСТ}---Искусство/Ремесло (Фотография), История, Использование Библиотек, Родной Язык, один навык для общения (Обаяние, Увещевание, Запугивание или Убеждение), Психология, два любых других навыка.
  \bigbreak
  \textbf{СЛЕДОВАТЕЛЬ}---Искусство/Ремесло (Лицедейство) или Маскировка, Огнестрельное Оружие, Юриспруденция, Слух, один навык для общения (Обаяние, Увещевание, Запугивание или Убеждение), Психология, Психология, Наблюдательность, один любой навык.
  \bigbreak
  \textbf{СЫЩИК}---Искусство/Ремесло (Фотография), Маскировка, Юриспруденция, Использование Библиотек, один навык для общения (Обаяние, Увещевание, Запугивание или Убеждение), Психология, Наблюдательность, один любой навык.
  \bigbreak
  \textbf{ПРОФЕССОР}---Использование Библиотек, Язык (Иной), Родной Язык, Психология, четыре любых других навыка как научные или личные специализации.
  \includegraphics[width=\linewidth]{img/bottom.png}
\end{cocpaperbox}

\subsection*{Шаг четвертый: предыстория}

Взгляните на навыки, характеристики и их значения. Вас захватывает воображение и вы начинаете представлять себе вашего персонажа во плоти. По мере того, как он вырисовывается все отчетливее, вы можете начать описывать его на бумаге. Кто он на самом деле? Где он вырос? Какая у него семья? Чем больше вы думаете о своем персонаже, тем более проработанным он окажется в итоге и тем интереснее будет играть за него в Зове Ктулху.

Каждый пункт жизнеописания (на обратной стороне листа персонажа) должен содержать небольшое краткое утверждение. Подумайте о том, каким персонаж будет казаться людям, встретившим его впервые и запишите ваши мысли в Личном Описании. Во что ваш персонаж верит и как относится к жизни? Опишите одним предложением его Идеологию/Верования. Какие маннеризмы ему присущи? Запишите его неприятные привычки и тому подобные вещи в графе Особенности. Вам не обязательно заполнять все графы в листе - достаточно начать с двух-трёх пунктов, чтобы персонаж не был статистом.
\smallbreak
\noindent \example{Полина записывает в графе Значимые Места ``Родилась и выросла в Нью-Йорке'', в графе Важные Вещи ``Не расстаюсь со своим верным револьвером'', а в графе Идеология/Верования - ``Наука может всему найти объяснение''.}

\subsection*{Шаг пятый: последние штрихи}

У вас уже есть кто-то, напоминающий завершенного персонажа. Взгляните в верхнюю часть листа персонажа и убедитесь, что вы подобрали ему имя, пол и возраст. На обратной стороне листа запишите оборудование, которое может быть у персонажа вашей Профессии.
\smallbreak
\noindent \example{Полина записывает ``Записная книжка, карандаши, перьевая ручка, расческа, заколки'', в графе Оборудование и Вещи, т.к. это то, что вполне может быть при себе у журналиста. Заколки вполне могут пригодиться, когда предстоит открывать замок!}
\smallbreak
На листе также есть пространство, куда вы можете вклеить портрет вашего исследователя. Если вы используете загружаемый PDF, то на слот под портрет можно просто кликнуть и загрузить изображение с компьютера.

Пока не беспокойтесь о графе Наличные и Имущество - она предназначена для продвинутой игры, где сбережения и наличка могут играть решающую роль (подробнее об этом можно прочитать в основной книге правил Зова Ктулху).

\section*{Игровая система}

Во время драматического накала, игра может потребовать ``броска на навык''. Проход по хорошо освещенному коридору не является драматической ситуацией, а вот бег по заваленному всяким мусором проходу, в то время как вас преследует толпа чудовищ - является!

Когда вам нужно совершить бросок на навык, вам нужно договориться с Хранителем о цели броска. Если бросок успешен, вы достигаете заявленной цели. Кроме того, если вы совершаете успешный бросок на навык, вам нужно отметить галочкой этот навык на листе персонажа. Подобную галочку можно поставить у каждого навыка только один раз. В конце сценария, этот навык может увеличиться, т.к. вы получили опыт в его использовании. Подробнее см. Награда за Успех на стр. 22.

Иногда может приходиться совершить бросок, для которого на листе персонажа нет навыка. В таком случае, вам нужно взглянуть на характеристики исследователя и определить, какая из них лучше подходит для данного случая и использовать ее как навык.

\subsection*{Броски на навык и уровни сложности}

Ваш Хранитель скажет вам, когда вам предстоит бросок на навык и то, насколько сложным он будет.
\smallbreak
\begin{itemize}[leftmargin=4mm]
  \item Несложная цель требует результата броска 1D100 равного или меньше вашего значения навыка (Обычный успех).
  \item Сложная цель требует результата равного или меньше половины вашего навыка (Тяжелый успех)
  \item Задача на пределе возможностей человека требует результата броска меньше или равной пятой значения вашего навыка (Экстремальный успех).
\end{itemize}
\smallbreak
Если вы можете обосновать действиями вашего исследователя такую необходимость - вы можете ``протолкнуть'' или ``пушнуть'' (от англ. ``to push'') ваш бросок. Пуш позволяет вам бросить кубики еще раз. Однако, в этот раз ставки выше.

% \begin{strip}
\begin{fullcocpaperbox}{}{}
  \raisebox{-10mm}{\includegraphics[width=\linewidth]{img/top.png}}
  \cocpaperboxtitle{Уровни успеха}
  (худший) ПРОВАЛ --- ОБЫЧНЫЙ УСПЕХ --- ТЯЖЕЛЫЙ УСПЕХ --- ЭКСТРЕМАЛЬНЫЙ УСПЕХ
  \includegraphics[width=\linewidth]{img/bottom.png}
\end{fullcocpaperbox}
% \end{strip}

Если вас постигнет неудача и в этот раз, Хранитель иожет навлечь ужасные неприятности на голову вашего персонажа. Перед тем, как вы пушнёте вам бросок, Хранитель может предсказать вам то, что случится, если вы провалите второй бросок. После этого игрок может оценить - хочет он пушить или нет.
\smallbreak
\noindent \example{ваш исследователь хочет приподнять тяжелую каменную дверь, преграждающую вход в склеп. Хранитель определяет сложность задачи как Тяжелую и просит вас сделать бросок на СИЛ, сообщая, что требуется ``Тяжелый успех''. Показатель СИЛ вашего персонажа равен 60, так что вам нужно выбросить 30 или меньше. Вы бросаете кубики, но выпадает 43---вы провалили бросок, т.к. выкинули больше половины вашего значения СИЛ. Вы спрашиваете, можно ли пушнуть бросок, говоря, что ваш персонаж удваивает усилия и применяет лопату в качестве рычага. Хранитель позволяет вам совершить второй бросок, но предупреждает, что если вас опять постигнет неудача, что вам не только не удастся открыть дверь, но и ``нечто'' может вас услышать и отправиться по вашему запаху!}

\subsection*{Противостоящие броски на навык}

Если два исследователя противостоят друг другу или один персонаж находится в противостоянии со значительным персонажем под контролем Хранителя, у которого есть собственные значения навыков, Хранитель может попросить совершить ``противостоящий бросок''.

Для того, чтобы его совершить, обе стороны кидают бросок на навык и сравнивают уровни своего успеха. Таким образом, Обычный успех бьет Неудачу, Тяжелый успех бьет Обычный, а Экстремальный - бьет Тяжелый. В случае ничьей побеждает тот, чье значение навыка выше. Если и значения навыка одинаковые - то решает бросок 1D100. Чье значение ниже - тот и выиграл.

\subsubsection*{Кубик Преимущества и Кубик Помехи}

Иногда окружение, условия и/или доступное время могут как помочь, так и помешать броску. При определенных условиях, Хранитель может предоставить ``преимущество'' либо ``помеху'' к броску. Каждая помеха вычитает одно преимущество. Преимущества и помехи действуют подобно увеличенной сложности и могут использоваться вместе с ней или вместо нее. Но, как правило, их используют для противостоящих бросков.
\smallbreak
\smallbreak
\noindent\textbf{За каждый кубик преимущества:} киньте дополнительный 1D10 отвечающий за десятки, а не единицы вместе с кубиками, используемыми в броске. То есть вы должны кинуть три кубика: один с единицами и два с десятками. Чтобы использовать преимущество, выберите лучший показатель на кубиках с десятками (более низкий).
\smallbreak
\noindent \example{два соперничающих исследователя, Малькольм и Хью борются за внимание Леди Грин. Только один из них может рассчитывать на ее благосклонность и руку, так что Хранитель определяет, что для понимания того, кто оказывается более успешным ухажером требуется противостоящий бросок. Игроки и Хранитель договариваются, что наиболее подходящим будет бросок на Обаяние. Хранитель вспоминает все, что произошло в ходе сценария к этому моменту: Малькольм посетил Леди Грин дважды, не забывая о дорогих подарках, в то время как Хью посетил ее лишь единожды и не позаботился о внешнем виде. Хранитель определяет, что Малькольм имеет преимущество и может использовать его в броске.

Игрок, играющий за Хью сперва кидает обычный бросок на навык. Значение его Обаяния равно 55, и игрок выбрасывет 40 на десятках и 5 на единицах, всего 45 - Обычный успех.

Игрок Малькольма также бросает обычный бросок на навык, но с преимуществом. Его Обаяние равно 50. Он кидает три кубика: два на десятки и один на единицы. Кубик с единицами показывает значение 4, а на десятках выпало 20 и 40. Не долго думая, игрок выбирает 20, в сумме 24 - Тяжелый успех.

Малькольм побеждает в противостоянии и его сватовство к Леди Грин принято благосклонно.
}
\smallbreak
\smallbreak
\noindent\textbf{За каждый кубик помехи:} киньте дополнительный 1D10 отвечающий за десятки, со стандартной парой кубиков. То есть вы должны кинуть три кубика: один с единицами и два с десятками. Чтобы получить помеху, вы должны худший показатель на кубиках с десятками (более высокий).
\smallbreak
\noindent \example{из-за череды неудачных стечений обстоятельств, два исследователя - Феликс и Гаррисон - оказались в плену у членов культа Алой Улыбки. Культисты решили ``поразвлечься'' с исследователями и хотят, чтобы они прошли Ордалию Боли, в ходе которой выживет только один. Проигравший будет принесен в жертву омерзительному божеству культа.

Ордалия Боли заключается в том, что нужно как можно дольше продержать тяжелый камень. Это требует противостоящего броска на СИЛ от обоих исследователей. Однако Хранитель определяет, что Гаррисон будет совершать бросок с помехой, ведь он недавно получил тяжелую травму (как раз тогда, когда был захвачен культистами) и все еще не восстановился.

Игрок Феликса выкидывает 51 на своем броске по СИЛ. Его показатель СИЛ равен 60---у него Обычный успех.

Сила Гаррисона - 55. Его игрок выкидывает 20 и 40 на кубиках с десятками и 1 на кубике с единицами. Он должен использовать более высокий результат из-за помехи---в итоге у Гаррисона тоже Обычный успех.

У обоих игроков Обычный успех, но Феликс выигрывает, т.к. его показатель характеристики СИЛ выше. Феликс удержал камень над головой дольше, чем изнемогающий от боли Гаррисон. Культисты визжат от удовольствия и тащат обессиленного Гаррисона к алтарю...
}

\subsubsection*{Броски Удачи}

Хранитель может просить совершить бросок на Удачу когда нужно определить обстоятельства вне контроля исследователей и отыграть слепую руку судьбы. Если, например, исследователь хочет узнать, есть ли рядом предмет, который можно использовать в качестве оружия или остался ли еще заряд в только что найденном фонарике - тогда нужен бросок на Удачу. Обратите внимание - если в ситуации уместнее кинуть навык или характеристику - тогда нужно кидать на них, а не на Удачу. Чтобы бросок на Удачу считался успешным, нудно выкинуть больше либо равно текущему показателю Удачи исследователя.

Если Хранитель просит совершить ``групповой бросок на Удачу'', это означает, что исследователь с худшим (самым низким) показателем Удачи из участвующих в сцене должен бросить кости от лица всей группы.
\smallbreak
\noindent \example{найти кэб - не требует броска кубиков, но найти кэб до того, как подозреваемые скроются на своей машине - уже зависит от броска Удачи. Показатель навыка Благосостояния может сыграть свою роль в привлечении внимания таксиста, который ищет щедрого на чаевые клиента. Но быстро найти такси в два часа утра в нехорошем квартале города уже скорее будет зависеть от Удачи. А вообще, найдется ли кэб? Нет навыка, который отвечал бы за местонахождения кэбов - это целиком воля судьбы, так что требуется бросок Удачи.}

\section*{Рассудок (РАС)}

Когда исследователь сталкивается с необъяснимыми ужасами Мифологии Ктулху или с чем-то естественным, но устрашающим (таким как вид обезображенного трупа друга), нужно прокинуть 1D100 по вашему текущему показателю рассудка. Если вы выбросите больше - вы потеряете большее число очков Рассудка, если меньше - либо вообще не потеряете, либо потеряете меньше. Потеря очков Рассудка обычно описывается в готовых сценариях как ``0/1D6'' или ``2/1D10''. Число слева от слеша говорит о том, сколько вы потеряете в случае успеха, а справа - сколько потеряете в случае неудачи.

\begin{fullcocpaperbox}{}{}
  \raisebox{-10mm}{\includegraphics[width=\linewidth]{img/top.png}}
  \cocpaperboxtitle{Примеры фобий и маний}
  Существуют сотни возможных фобий и маний. Здесь приведены лишь некоторые из возможных.
  \begin{multicols}{2}
  \subsubsection*{Фобии}
  \begin{itemize}[leftmargin=4mm]
  \item Боязнь высоты (акрофобия).
  \item Боязнь пауков (арахнофобия).
  \item Боязнь книг (библиофобия).
  \item Боязнь зеркал (эйзоптрофобия).
  \item Боязнь крови (гемафобия).
  \item Боязнь мертвых существ (некрофобия).
  \item Боязнь зубов (одонтофобия).
  \item Боязнь огня (пирофобия).
  \item Боязнь телефонов (телефонофобия).
  \item Боязнь незнакомцев или чужаков (ксенофобия).
\end{itemize}
  \bigbreak
  \bigbreak
  \bigbreak  \bigbreak
  \bigbreak
  \bigbreak
  \subsubsection*{Мании}
  \begin{itemize}[leftmargin=4mm]
  \item Патологическая вежливость (агатомания).
  \item Одержимость болью (алгомания).
  \item Иррациональная жизнерадостность (аменомания).
  \item Одержимость кражей книг (библиоклептомания).
  \item Одержимость вершением справедливости (дикемания).
  \item Неконтролируемое желание смеяться (гелиомания).
  \item Иррациональное желание кричать (клазомания).
  \item Иррациональная потребность в кражах (клептомания).
  \item Уверенность в наличии воображаемой болезни (носомания).
  \item Иррациональное желание лгать (псевдомания).
\end{itemize}
  \end{multicols}
  \includegraphics[width=\linewidth]{img/bottom.png}
\end{fullcocpaperbox}

Когда вы проваливаете бросок на Рассудок, Хранитель отыгрывает ваше следующее действие самостоятельно, ведь вас охватывает неконтролируемый страх: возможно, вы неожиданно вскрикиваете или невольно нажимаете на курок вашего револьвера.

Если исследователь теряет 5 или больше очков Рассудка в ходе одного броска на Рассудок, это означает, что они получают серьезную эмоциональную травму. Игрок должен кинуть 1D100. Если результат меньше или равен показателю Интеллекта (ИНТ) его персонажа, это означает, что персонаж вполне понял, что перед ним предстало и впадает во временное безумие (длительностью в  1D10 часов). Если игрок провалил бросок, то их разум инстинктивно отказался воспринимать ужас и они остаются (пока что) в здравом уме.

В дополнение, безумный исследователь страдает от ``приступа безумия'' --- киньте 1D10 и обратитесь к Таблице Приступов Безумия (стр. 17). Если исследователь находится в компании своих товарищей, то отыгрываете эффект от раунда к раунду. Если исследователь один, то вы можете использовать этот результат, чтобы описать, как его находят через некоторое время в плохом состоянии, возможно запертым изнутри в шкафу, или пьяным вдрызг в канаве.
% \begin{strip}
\begin{fullcocpaperbox}{}{}
  \raisebox{-10mm}{\includegraphics[width=\linewidth]{img/top.png}}
  \cocpaperboxtitle{Таблица Приступов Безумия}
  Хранитель выбирает самостоятельно или путем броска 1D10.
  \begin{multicols}{2}
  \begin{enumerate}
    \item \textbf{Амнезия:} исследователь не помнит о событиях, которые произошли с момента их последнего пребывания в безопасном месте.
    \item \textbf{Психосоматическая инвалидность:} исследователь испытывает психосоматическую слепоту, глухоту или способность управлять конечностью на 1D10 раундов.
    \item \textbf{Беспричинная жестокость:} кровавая пелена застилает глаза исследователя и под воздействием неконтролируемой ярости персонаж начинает набрасываться на все, что оказывается под рукой, на союзников и врагов в том числе. Эффект длится 1D10 раундов.
    \item \textbf{Паранойя:} Исследователь испытывает сильную паранойю на протяжении 1D10 раундов. Все хотят ему навредить! Никому нельзя доверять! За ним следят, его предали, все, что он видит - это обман.
    \item \textbf{Важное лицо:} посмотрите в список Важных Лиц на листе персонажа. Исследователь принимает другого человека, присутствующего в сцене, за одно из Важных Лиц. Исследователь действует в соответствии со своим отношением к нему. Эффект длится 1D10 раундов.
    \item \textbf{Обморок:} персонаж теряет сознание, приходя в себя через 1D10 раундов.
    \item \textbf{Паническое бегство:} исследователя охватывает неодолимое желание бежать куда глаза глядят, как можно дальше от этого места, даже если ему придется использовать единственное транспортное средство, оставив товарищей на произвол судьбы.
    \item \textbf{Истерика:} исследователь охвачен смехом, его душат рыдания, он кричит во всю мощь легких и т.п. на протяжении 1D10 раундов.
    \item \textbf{Фобия:} персонаж приобретает новую фобию, например клаустрофобию (боязнь замкнутых пространств), демонофобию (боязнь духов / демонов) или катсаридофобию (боязнь тараканов). Даже если источник фобии не присутствует, исследователь думает, что это не так на протяжении 1D10 раундов, а все действия совершаются с помехой на протяжении действия эффекта.
    \item \textbf{Мания:} исследователь приобретает новую манию, например аблютоманию (неодолимое желание вымыться), псевдоманию (тягу ко лжи) или гельминтоманию (одержимость любовью к червям). Исследователь стремиться удовлетворить манию на протяжении ближайших 1D10 раундов, а все действия приобретают помеху, пока длится эффект.
  \end{enumerate}
  \end{multicols}
  \includegraphics[width=\linewidth]{img/bottom.png}
  \end{fullcocpaperbox}{}{}
% \end{strip}
Если ваш исследователь временно безумен, Хранитель может решить добавить фобию или манию на вашем листе персонажа (такие как ``боязнь темноты'', ``боязнь замкнутых пространств'', ``клептомания''). Как альтернатива, Хранитель также может переписать одну из граф в вашей предыстории, искажая ее (там, где вы написали ``доверчивая'' в качестве Черты, Хранитель может переписать на ``боязливая'').

Пока исследователь находится в состоянии временного безумия, Хранитель может наградить его ``Бредом'' (галлюцинациями), например персонаж не может понять - за его ногу цепляется зомби или это лишь бездомный человек попросил мелочь? В таком случае игрок может быть уверен только если он прокинет ``определение действительности'', для чего нужно опять совершить бросок на Рассудок. Если он успешен, персонаж может понять, что реально, а что нет. Если не успешен - персонаж погружается еще глубже в пучину безумия!

После прохождения 1D10 часов, исследователь приходит в себя и не может быть подвержен галлюцинациям, но измененная предыстория, фобии и мании остаются.
  \begin{fullcocpaperbox}{}{}
  \includegraphics[width=\linewidth]{img/top.png}
  \cocpaperboxtitle{Памятка Рассудка}
  \begin{multicols}{2}
    Когда совершается бросок на Рассудок, обычно потеря отображается в таком виде: XX/XX (например, 1/1D6). Число слева - это количество очков Рассудка, которое будет потеряно в случае успешного броска, справа - в случае неудачи.

    \textbf{Непроизвольное действие:} любая потеря очков Рассудка приведет к непроизвольному действию, определяемому Хранителем (невольный вскрик, оборок на один раунд, случайное нажатие на спусковой крючок).

    \textbf{Если в ходе броска было потеряно больше 5 очков Рассудка:} попросите исследователя совершить бросок на Интеллект (ИНТ): если его постиг провал, персонаж остается в здравом рассудке, если он был успешен (т.е. выкинул значение равное или меньшее своему показателю ИНТ), тогда исследователь временно теряет рассудок (на 1D10 раундов). Хранитель может применить нижеследующее, как ему самому больше нравится:

    \begin{enumerate}
    \item Приступ безумия: выберите или бросьте на один из вариантов в Таблице Приступов Безумия и примените результат на 1D10 раундов.
    \item (Опционально): возьмите лист персонажа и добавьте ему подходящую случаю запись в предысторию или измените существующую, основываясь на виде безумия или его причине.
    \item  (Опционально): добавьте к предыстории персонажа фобию или манию.
    \item  Опишите галлюцинации, которые может испытывать персонаж, в случае если он не сможет прокинуть еще один бросок на Рассудок. Галлюцинации могут преследовать персонажа не более чем на 1D10 часов.
    \end{enumerate}
  \end{multicols}
  \includegraphics[width=\linewidth]{img/bottom.png}
\end{fullcocpaperbox}

\section*{Бой}

Когда вам противостоят ужасы Мифологии Ктулху, чаще всего верным решением будет развернуться и убежать. Еще лучше - изначально избегать конфронтации, потому что большинство чудовищ крайне могущественны и обычной пулей их не повредить. Тем не менее, иногда нет иного выхода, кроме как вломиться к культистам с оружием наперевес.

Когда происходит бой, все персонажи, как исследователи, так и их оппоненты, ходят в порядке показателя ЛВК. Чем выше этот показатель, тем раньше действует персонаж.

\smallbreak
\noindent \example{Персонажу Лёши противостоит культист, только что призвавший чудовище. Кажется, назревают неприятности. У исследователя показатель ЛВК равен 50, культиста 45, чудовища 70. Таким образом, чудовище в рамках раунда боя будет ходить первым, персонаж Алексея - вторым, культист - третьим.}

Продолжительность раунда боя в Зове Ктулху может быть описана как ``достаточная, чтобы все совершили одно действие''. Хранитель контролирует течение боя. В ход каждого персонажа Хранитель решает (или спрашивает у игрока, если ходит не NPC) что он будет делать. Чаще всего это нечто простое, вроде "я нападаю на чудище", "я достаю револьвер" или "я убегаю". Хранитель должен предоставить всем возможность сделать что-то быстрое, стараясь не замедлять ритм повествования.

У исследователей есть три боевых навыка: "Бой", "Уклонение" и и "Огнестрельное Оружие". Два из этих навыков состоят из нескольких специализаций, например Бой (Рукопашный) или Огнестрельное Оружие (Винтовка / Дробовик). Каждый игрок должен решить, имеет ли их персонаж специализацию и если имеет, то какую, в момент распределения очков по навыкам. Обратите внимание, что под специализацию Бой (Рукопашный) подпадает бой без оружия и использование простого вооружения, вроде ножей и дубин, однако более специализированное оружие типа меча потребует специализации Бой (Меч).

Вы совершаете броски в бою, используя наиболее подходящие навыки, как и с любыми другими бросками. Единственное исключение - их нельзя пушить.

\subsection*{Ближний Бой}

В ход персонажа в порядке  показателя ЛВК он может инициировать атаку против оппонента. Кроме того, каждый раз, когда на персонажа нападают, он может выбрать как они отреагируют: либо уклонением (пытаясь полностью избежать урона), либо контратакуя (пытаясь избежать, блокировать или парировать атаку, в то же самое время стараясь дать сдачи).

И нападающий и защищающийся кидают процентный кубик (1D100) и сравнивают результаты бросков:

\begin{itemize}[leftmargin=4mm]
  \item Если вы контратакуете, то вы используете своей навык Боя. Вам нужно достичь большего уровня успеха, чем у оппонента.
  \item Если вы уклоняетесь, вы используете навык Уклонения. Атакующий вас персонаж должен достичь большего уровня успеха, чем вы.
\end{itemize}

Все довольно просто: выигравшая сторона не получает никакого урона и нанесет урон (если она не уклонялась) своему противнику.

Когда персонаж контратакует, самый большой уровень урона, который он может нанести - это ``обычный'', в то время как персонаж, инициирующий атаку может достичь максимум ``экстремального'' урона (см. ниже).

\smallbreak
\noindent \example{Вурдалак пытается попасть своей когтистой лапой по персонажу Полины, Полина хочет уклониться. Хранитель выкидывает 03 - Экстремальный успех (значение меньше пятой показателя Боя вурдалака). Полина выкидывает 20 на своем броске Уклонения - Сложный успех. Нападающий достиг более высокого уровня успеха, чем уклоняющийся, поэтому по исследователю Полины попали, и он получает максимум урона - 10 (1D6 + 1D4), потому что вурдалак достиг Экстремального успеха.

Вурдалак - чудовище, которое может выполнить 3 атаки за раунд (все они происходят одновременно). Своей следующей атакой он попытается укусить персонажа Полины, которая хочет контратаковать. Полина достигает Сложного успеха, вурдалак - Обычного. Полина преуспела больше, чем оппонент, поэтому она не только избегает еще одних повреждений, но и наносит 1D3 урона врагу.

Помните: если результат Уклонения равен результату нападающего, выигрывает уклоняющийся. Если аналогичная ситуация с контратакой - то побеждает атакующий, а не контратакующий. Уклоняться проще, чем контратаковать.
}
\smallbreak

\subsubsection*{Экстремальный Урон}

Атаки, которые достигли Экстремального уровня успеха наносят увеличенный урон.

\begin{itemize}[leftmargin=4mm]
  \item Тупое оружие наносит максимальный урон плюс максимальный бонус к урону (если он есть).
  \item Проникающее оружие (клинки и пули) наносит максимальный урон плюс максимальный бонус к урону (если он есть) плюс дополнительный бросок на урон оружием (1D10 + 10 урона для пистолета, например).
\end{itemize}

\smallbreak
\noindent \example{Персонаж Ромы выигрывает раунд боя с Экстремальным успехом. У него в руках дубина (тупое оружие) и его бонус к урону равен 1D4. Атака наносит 6 + 4 = 10 урона. Если бы он использовал нож (проникающее оружие), то нанес бы 4 + 4 + 1D4 урона.
}
\smallbreak

\subsection*{Правила для огнестрельного оружия}

Стреляющий из огнестрельного оружия кидает процентный кубик и сравнивает результат со своим показателем навыка Огнестрельное Оружие.

\begin{itemize}[leftmargin=4mm]
  \item Если оружие было взято наизготовку, то при определении очередности хода это дает бонус в +50 к показателю ЛВК.
  \item Если вы совершаете в один раунд 2 или 3 выстрела из пистолета, вы должны добавить один кубик помехи к каждому выстрелу.
  \item При выстреле в упор (выстрел считается таковым, если расстояние до цели меньшей 1/5 ЛВК в футах (0.3 метра), стреляющий получает кубик преимущества на бросок.)
\end{itemize}

Цель выстрела не может выбрать реакцию (невозможно увернуться от пули), но можно ``броситься в укрытие'', сделав бросок на Уклонение. Если бросок успешен, то атакующий совершает броски на попадание с кубиком помехи. Персонаж, который решает укрыться теряет возможность совершить свою следующую атаку (независимо от того, успешно ли он укрылся). Если персонаж уже совершил свою атаку в этом раунде, то он теряет возможность совершить атаку в следующем.

\smallbreak
\noindent \example{Исследователь Лёши взял револьвер наизготовку, а затем увидел, что к нему, размахивая ятаганом несется культист. ЛВК персонажа равна 50, но его взведенное оружие дает бонус в 50, давая показатель в 100 при определении очередности хода. ЛВК культиста всего 45, так что исследователь Алексея будет ходить первым.

Культист видит револьвер и бросается в укрытие (совершая бросок на Уклонение) и добивается успеха. Лёша совершает свой бросок с помехой, в результате проваливая его. Культист потерял своё действие в раунде, т.к. бросился в укрытие, поэтому раунд заканчивается. Новый раунд начинается, давая еще одну возможность Алексею прострелить культисту голову.
}

\begin{cocpaperbox}{}{}
  \raisebox{-10mm}{\includegraphics[width=\linewidth]{img/top.png}}
  \cocpaperboxtitle{Оружие и урон}
  \bigbreak
\begin{itemize}[leftmargin=4mm]
  \item Безоружные атаки (человек): 1D3 + бонус к урону
  \item Небольшой нож: 1D4 + бонус к урону
  \item Мачете: 1D8 + бонус к урону
  \item Небольшая дубина: 1D6 + бонус к урону
  \item Бейсбольная бита: 1D8 + бонус к урону
  \item Пистолет: 1D10
  \item Дробовик: 4D6 (на короткой дистанции*, в противном случае 2D6, без проникновения)
  \item Винтовка: 2D6 + 4
\end{itemize}
\bigbreak
*Короткая дистанция: в пределах показателя ЛВК в футах (т.е. если ЛВК равна 60, то короткой дистанцией считается расстояние в 60 футов (18 метров))
  \includegraphics[width=\linewidth]{img/bottom.png}
\end{cocpaperbox}

\smallbreak

\subsection*{Маневры в бою}

Если игрок заявляет цель отличную от обычного нанесения физического урона, подобная цель может быть описана путем ``Боевого Маневра''. Успешный маневр позволяет игроку достичь одной цели, например:

\begin{itemize}[leftmargin=4mm]
  \item Разоружить противника
  \item Сбить противника с ног
  \item Схватить и удерживать врага, из-за чего противник должен совершать все броски с помехой, пока не избавится от захвата
\end{itemize}

Маневр воспринимается как обычная атака, бросок совершается по навыку Бой (Рукопашный). Оппонент может выбрать уклонение или контратаку как обычно. Затем сравнивается Комплекция. Если персонаж, который совершает маневр, имеет меньший показатель Комплекции, то он совершает маневр с дополнительным кубиком помехи за каждое очко разницы (до максимума в два кубика помехи). Если персонаж имеет показатель Комплекции, меньший Комплекции противника на 3 или больше, то маневр бесполезен: персонаж может схватить оппонента, но у него не хватит силы, чтобы сдвинуть его с места.

\smallbreak
\noindent \example{Полина пытается вытолкнуть вурдалака из ближайшего окна (рассматривается как маневр). Комплекция персонажа Полины равна 0, Комплекция вурдалака 1, поэтому Полина совершает бросок с помехой. Полина выкидывает 02 и 22, из-за помехи она вынуждена выбрать худшее (наивысшее) значение. У нее - Тяжелый успех (22 меньше 1/2 показателя навыка Боя Полины). Вурдалак совершает контратаку и выбрасывает Обычный успех на навыке Боя. Полина достигла более высокого уровня успеха и с жутким воплем вурдалак летит к земле.
}

\subsubsection*{Окружены}

Персонаж, которого окружили враги оказывается в трудной ситуации. После того, как персонаж совершил уклонение или контратаку в текущем раунде, все последующие атаки по нему совершаются с преимуществом. Это правило не применяется к огнестрельному оружию.

\smallbreak
\noindent \example{Вурдалак имеет 3 атаки, у исследователя Полины - 1. После первой атаки вурдалака Полина может выбрать, уклоняться или контратаковать как обычно. Но во время второй и третьей атаки вурдалака, он получает преимущество, т.е. по сути Полина расценивается как окруженная.

В другой ситуации, Рома противостоит двум культистам, он тоже окружен. Атака первого культиста будет совершаться как обычно, второй культист уже получит преимущество.
}

\section*{Очки здоровья, раны и исцеление}

Очки урона вычитаются из очков здоровья персонажа. Когда очки здоровья падают до нуля, исследователь теряет сознание и в некоторых случаях может умереть.

Когда персонаж получает за один удар урон равный или больший чем половина его полного значения очков здоровья, он получает ``значительную рану'' - ему нужно прокинуть его значение ТЛС или он теряет сознание. Если исследователь со значительной раной достигает нуля очков здоровья, он начинает умирать. В конце каждого последующего раунда включая этот ему нужно пробросить бросок на ТЛС или умереть. Только успешное применение навыка Первой Помощи может снять это состояние, стабилизируя персонажа. Если персонаж получает за один удар урон больший или равный его полному показателю очков здоровья - он умирает моментально.

Обратите внимание, что если у персонажа нуль очков здоровья, но нет значительной раны, то он не умирает: смерть возможна только при наличии значительной раны.

\begin{itemize}[leftmargin=4mm]
  \item Персонажи без значительной раны исцеляют 1 очко урона в день естественным путем.
  \item Персонажи со значительной раной должны совершить бросок на исцеление (выбросив значение меньшее или равное их показателю ТЛС) в конце каждой недели. При успехе, они восстанавливают 1D3 очков здоровья (2D3 в случае Экстремального успеха). Значительная рана исчезает, если персонаж выбросли на броске исцеления Экстремальный успех или текущее значение очков здоровья достигло половины или более их максимального значения. Таким образом, для полного исцеления могут потребоваться недели.
\end{itemize}

Успешное применение навыка Первой Помощи также может восстановить 1 очко здоровья или привести персонажа в чувство. Если персонаж был при смерти, Первая Помощь может его стабилизировать, а затем при помощи навыка Медицины персонажу можно восстановить 1D3 очков здоровья, но для этого так же потребуется минимум один час и необходимые инструменты и лекарства. Если навык Медицины используется на умирающем персонаже, в конце недели он сможет совершить бросок на исцеление.

\smallbreak
\noindent \example{Изначально у персонажа Лёши 12 очков здоровья. В понедельник в баре он оказывается участником драки и получает урон от трех ударов в челюсть, стоивших ему 4, 2 и 4 очка здоровья. Всего это 10 очков здоровья, у исследователя Алексея их остается 2. Он не получил значительной раны (ни одна атака не нанесла большого количества урона), так что он будет восстанавливаться со скоростью 1 очко здоровья в день. В четверг персонаж (теперь с 5 очками здоровья) по неловкости выпадает из ока, получая 7 очков урона. Это - значительная рана (7 больше половины максимального количества очков здоровья персонажа). Друг применяет Первую Помощь и спешно отвозит его в больницу. Посл того, как прошло семь дней, Алексей успешно прокидывает ТЛС и на броске 1D3 он выкидывает 2 и восстанавливает 2 очка здоровья. В конце второй недели Алексей выбрасывает Экстремальный успех на броске ТЛС и получает 4 очка здоровья на броске 2D3, его текущее количество здоровья равно 6. Это убирает у него значительную рану (он восстановил половину своего максимального значения очков здоровья), после чего он исцеляется со скоростью 1 очков здоровья в день.
}

\subsection*{Прочие виды урона}

Часто Хранителю приходится определять количество урона, которое персонажи получают от разнообразных случайных происшествий. Независимо от причины, определитесь с возможной травмой и оцените ее по левой колонке в Таблице Прочих Видов Урона. Каждая травма - от каждого происшествия или одного раунда боя (одна травма от одного раунда, в котором персонажа били и попали пулей, одна травма от другого раунда в котором персонажа топили, одна - от раунда в котором он горел и т.д.).

\begin{fullcocpaperbox}{}{}
  \raisebox{-10mm}{\includegraphics[width=\linewidth]{img/top.png}}
  \cocpaperboxtitle{Таблица Прочих Видов Урона}
  \rowcolors{2}{}{TableRow}
  \begin{tabularx}{170mm}{XYX}
  \rowcolor{TableHeading}\color{CoCPaperBox}\bfseries Травма & \color{CoCPaperBox}\bfseries Урон & \color{CoCPaperBox}\bfseries Примеры \\
  \textbf{Малая:} человек может выжить, испытав такую урон несколько раз. & 1D3 & Удар рукой, ногой, головой, слабая кислота, нахождение в задымленном помещении*, удар камнем средних размеров, падение (урон рассчитывается за каждые 10 футов (0.3 метра) падения) в мягкую болотную почву \\
  \textbf{Средняя:} может привести к серьезной ране, убить можно меньшим количеством подобных атак. & 1D6 & Падение (за 10 футов) на траву, удар дубиной, сильная кислота, дыхание под водой*, нахождение в вакууме*, малокалиберная пуля, стрела, огонь (прикосновение к огню). \\
  \textbf{Тяжелая:} скорее всего приведет к тяжелой ране. Одна или две могут привести к смерти. & 1D10 & Пуля калибра .38, падение (за 10 футов) на бетон, удар топором, огонь (атака из огнемета, бег через горящую комнату), нахождение в 6-9 метрах от разорвавшейся гранаты или динамитной шашки, слабый яд**.\\
  \textbf{Смертельная:} Среднестатистический человек умрет в 50\% случаев. & 2D10 & Быть сбитым машиной на скорости 50 км/ч, нахождение в 3-6 метрах от разорвавшейся гранаты или динамитной шашки, сильный яд**. \\
  \textbf{Терминальная:} Вероятна моментальная смерть. & +4D10 & Быть сбитым машиной на полной скорости, нахождение в 3 метрах от разорвавшейся гранаты или динамитной шашки, смертельный яд**. \\
  \textbf{В лепешку:} Почти наверняка моментальная смерть. & +8D10 & Попадание в аварию на полной скорости, быть сбитым поездом. \\
  \end{tabularx}
\smallbreak
*Утопление и удушение: каждый раунд необходимо совершать бросок на ТЛС, если бросок провален, то каждый последующий раунд наносится урон до смерти или до того момента, как жертва сможет дышать. Смерть наступает при достижении нуля очков здоровья (правило значительной раны игнорируется).

**Яды: Экстремальный успех при броске ТЛС ополовинивает получаемый от яда урон.

  \includegraphics[width=\linewidth]{img/bottom.png}
\end{fullcocpaperbox}

\section*{Награда за успех}

Несмотря на то, что многие из сценариев к Зову Ктулху могут быть сыграны за один вечер, некоторые из них подразумевают, что исследователи будут расти и развиваться. В конце сценария или сессии, Хранитель должен попросить сделать ``бросок на увеличение навыков''. В этот момент каждый игрок кидает процентный кубик (1D100) за каждый навык, отмеченный галочкой (т.е. те, которые на протяжении сессии использовались успешно). Если этот бросок будет провален (т.е. значение будет выше текущего показателя навыка), это означает, что исследователь получил новый опыт в его использовании и вы можете добавить 1D10 очков к этому навыку. Другими словами - чем лучше вы владеете навыком, тем сложнее вам в этом навыку достичь новых высот.

\smallbreak
\noindent \example{Анна (исследователь Маши) совершает успешный бросок на Наблюдательность в течение игры, и Маша помечает этот навык галочкой на листе персонажа. После окончания сценария, Хранитель просит Машу кинуть на увеличение навыка. Уровень Наблюдательности Анны равен 45\%, а бросок равен 43. Увеличения навыка не происходит. Если бы Маша выкинула больше чем 45, она бы могла добавить 1D10 очков к своему показателю Наблюдательности. Маша стирает галочку с листа персонажа. В случае успеха в проверке Наблюдательности в следующей игре, она сможет отметить его снова.
}

\begin{fullcocpaperbox}{}{}
  \raisebox{-10mm}{\includegraphics[width=\linewidth]{img/top.png}}
  \cocpaperboxtitle{Описание навыков}
  \begin{multicols}{3}
  Большинство навыков означают ровно то, что написано в названии: например, Вождение Автомобиля, который определяет, насколько хорошо персонаж управляется с автомобилем, или Скалолазание, которое говорит о том, насколько хорошо исследователь сможет вскарабкаться по стене здания. Тем не менее, наименование некоторых навыков может ввести в заблуждение, поэтому здесь приведено краткое описание навыков, указанных на листе персонажа.
  \smallbreak
  \textbf{Обратите внимание:} на листе исследователя рядом с каждым навыком в скобках приведено процентное значение: это показатель навыка по умолчанию. То есть любой человек может использовать пистолет (базовое значение Огнестрельного Оружия равно 20\%), даже если он раньше никогда не стрелял.
  \smallbreak
  \textbf{Счетоводство:} понимание бухгалтерских процедур, может сказать о финансовом состоянии человека или компании.
  \smallbreak
  \textbf{Антропология:} определить и понять образ жизни человека (или народа) путем наблюдения.
  \smallbreak
  \textbf{Оценка:} узнать ценность конкретного предмета, включая то, насколько качественно он сделан, из каких материалов и в каком регионе.
  \smallbreak
  \textbf{Археология:} позволяет датировать и опознать артефакты ушедших культур, а также определить подделку.
  \smallbreak
  \textbf{Искусство и Ремесло:} владение этим навыком дает возможность создавать или ремонтировать предметы или сущности, которые могут быть произведением искусства (например картина или песня) или ремесла (столярного или кулинарного). Выберите подходящую специализацию и запишите в месте, предоставленном на листе персонажа.
  \smallbreak
  \textbf{Обаяние:} это скорее совокупность вещей: то, как персонаж выглядит, его умение соблазнять, льстить или просто производить приятное впечатление. Обаяние может быть использовано чтобы заставить кого-то вести себя определенным образом, но при этом не в манере полностью противоположной его обычному поведению. Этому навыку можно противостоять Обаянием или Психологией.
  \smallbreak
  \textbf{Скалолазание:} определяет не только способность карабкаться по скалам, сюда подпадают и деревья, стены и другие вертикальные поверхности, как с использованием снаряжения, так и без него.
  \smallbreak
  \textbf{Благосостояние:} уровень Благосостояния персонажа отвечает на вопросы ``насколько он богат?'' и ``к какому слою общества он принадлежит?''. Чем больше очков потрачено на этот навык, тем он богаче. В зависимости от того, сколько очков профессиональных навыков вы потратите, ваш персонаж будет:
    \begin{itemize}[leftmargin=4mm]
      \item Благосостояние 0: без гроша в кармане, живет на улице.
      \item Благосостояние 1-9: бедняк, владеет самым минимумом.
      \item Благосостояние 10-49: среднего достатка, может обеспечить себе сносный уровень жизненного комфорта.
      \item Благосостояние 50-89: зажиточный, может позволить себе некоторую степень роскоши.
      \item Благосостояние 90-98: богатый, огромные доходы и непозволительная для всех остальных роскошь.
      \item Благосостояние 99: сказочно богатый, деньги не имеют значения.
    \end{itemize}
  \smallbreak
    \noindent \example{Рома выбирает своему солдату уровень Благосостояния равный 40\%, это означает что у того будет средний уровень дохода.
    }
  \smallbreak
  \textbf{Мифология Ктулху:} отражает понимание нечеловеческой природы Мифологии Ктулху. Он не базируется на преумножении знаний, в отличие от научных навыков. Он скорее описывает открытие и подстраивание человеческого разума к истинам Мифологии. Таким образом, очки навыков для Мифологии приобретаются от личных встреч с Мифологией (чудовищами или омерзительными тайнами, описанными в редких книгах). Секреты Мифологии с трудом воспринмаются человеческим разумом и беспощадны к психике. Ни один исследователь не может вкладывать в этот навык очки (если это заранее не оговорено с Хранителем).
  \smallbreak
  \textbf{Маскировка:} используется, когда вы хотите выдать себя за кого-то еще.
  \end{multicols}
  \includegraphics[width=\linewidth]{img/bottom.png}
\end{fullcocpaperbox}

\begin{fullcocpaperbox}{}{}
  \raisebox{-10mm}{\includegraphics[width=\linewidth]{img/top.png}}
  \cocpaperboxtitle{Описание навыков}
  \begin{multicols}{3}
  \textbf{Уклонение:} позволяет исследователю инстинктивно избегать ударов, брошенных снарядов и т.д. Персонаж может попробовать уклониться любое количество раз за бой (но только один раз за атаку). Если атаку можно увидеть, то от нее можно попробовать уклониться, т.е. от пуль, двигающихся со скоростью не видной глазу, уклоняться нельзя. Лучшее, что может сделать персонаж - это попытаться броситься в укрытие. Определите стартовое значение, поделив пополам ваше Обычное значение ЛВК.
  \smallbreak
  \textbf{Вождение Автомобиля:} умение водить машину или грузовик, совершать стандартные маневры и устранять обычные неисправности. Если исследователь хочет сбросить преследователя с хвоста или наоборот, догнать чужую машину - бросок на Вождение будет вполне уместным.
  \smallbreak
  \textbf{Электрика:} починка и перенастройка электрического оборудования, таких как автомобильное зажигание, электрические моторы, детонаторы и сигнализации.
  \smallbreak
  \textbf{Увещевание:} этот навык целенаправленно ограничен вербальными уловками, обманом и пусканием пыли в глаза. Он подходит для того, чтобы ошарашить вышибалу на входе в ночной клуб, уговорить кого-либо не глядя подписать документ, заставить полисмена закрыть глаза на происшествие и так далее. Этому навыку можно противостоять с помощью Увещевания или Психологии. 
  \smallbreak
  \textbf{Бой:} умение персонажа драться в ближнем бою. Вы можете выбрать любую специализацию, такую как Рукопашный бой (сюда входят ножи и дубины, кулачный бой и боевые искусства), бой на Мечах, Топорах, Копьях или Кнутах.
  \smallbreak
  \textbf{Огнестрельное Оружие:} покрывает владение огнестрельным оружием, а также луками и арбалетами. Вы можете тратить очки навыков на любые специализации, включая Пистолеты, Винтовки / Дробовики, Луки или Арбалеты.
  \smallbreak
  \textbf{Первая Помощь:} экстренная медицинская помощь. Этот навык не может использоваться для лечения заболеваний (для этого требуется навык Медицина). Для того, чтобы применение возымело эффект, Первая Помощь должна быть применена не позднее часа с момента получения ранения, в таком случае восстанавливается 1 очко здоровья, и персонаж может придти в сознание.
  \smallbreak
  \textbf{История:} позволяет вспомнить деталь или событие из истории страны, города, региона или личности, в зависимости от обстоятельств.
  \smallbreak
  \textbf{Язык, Иной:} когда выбирается этот навык, нужно выбрать конкретный язык и указать его рядом с навыком. Персонаж может знать любое количество языков, но за каждый необходимо отдельно платить очками навыков. Этот навык отображает способность исследователя понимать, говорить, читать и писать на определенном неродном языке.
  \smallbreak
  \textbf{Язык, Родной:} вы выбираете любой язык, которым ваш персонаж владеет лучше всего. Стартовое значение равно обычному показателю ОБР.
  \smallbreak
  \textbf{Юриспруденция:} отображает вероятность знания относящегося к делу закона, прецедента, юридической уловки или судебной процедуры. Полезен при взаимодействии с полицией, юристами и судами.
  \smallbreak
  \textbf{Использование Библиотек:} умение находить информацию в библиотеках, такую как конкретная книга, газета, полезная ссылка (если в конкретной библиотеке подобная информация присутствует). Использование этого навыка подразумевает трату нескольких часов, потраченных на поиски.
  \end{multicols}
  \includegraphics[width=\linewidth]{img/bottom.png}
\end{fullcocpaperbox}

\begin{fullcocpaperbox}{}{}
  \raisebox{-10mm}{\includegraphics[width=\linewidth]{img/top.png}}
  \cocpaperboxtitle{Описание навыков}
  \begin{multicols}{3}
  \textbf{Использование Библиотек:} умение находить информацию в библиотеках, такую как конкретная книга, газета, полезная ссылка (если в конкретной библиотеке подобная информация присутствует). Использование этого навыка подразумевает трату нескольких часов, потраченных на поиски.
  \smallbreak
  \textbf{Слух:} интерпретация и понимание звуков. Сюда относится подслушивание разговоров, бормотания за закрытыми дверями и шепотов в кафе.
  \smallbreak
  \textbf{Слесарничество:} умение открывать автомобильные двери, замыкать зажигание, взламывать библиотечные окна, разгадывать головоломки на комбинаторику и обходить стандартные системы сигнализации. Персонаж, хорошо владеющий этим навыком может изготавливать ключи (при наличии образца), отмычки и прочие инструменты.
  \smallbreak
  \textbf{Механика:} отвечает за возможность починить или создать новый механизм. Персонаж может осуществлять базовые кровельные и сантехнические работы, создавать инженерные устройства, такие как блок или паровая помпа. Может использоваться для открытия стандартных домовых замков, но более сложных замков может потребоваться Слесарничество.
  \smallbreak
  \textbf{Медицина:} диагностировать и лечить последствия несчастных случаев, травм, болезней, отравлений и т.д. Лечение требует как минимум одного часа и может применяться в любое время после получения урона, но если навык не применен в тот же день, сложность броска увеличивается (см. стр. 14) и требует Сложного успеха. Персонаж, успешно полеченный Медициной восстанавливает 1D3 очка здоровья (в дополнение к полученным от Первой Помощи ), кроме случаев, когда персонаж при смерти, тогда до применения Медицины исследователь должен быть стабилизирован успешным броском Первой Помощи.
  \smallbreak
  \textbf{Природа:} отражает традиционное (ненаучное) знание и личные наблюдения фермеров, рыбаков, охотников и увлекающихся любителей. С его помощью можно идентифицировать растения, животных, в общих чертах их привычки и ареал обитания, определять их следы, запахи и звуки.
  \smallbreak
  \textbf{Навигация:} выбрать правильный путь к какой-либо точке, будь то в незнакомом городе или в глуши. Дает возможность эффективно читать карты и определять расстояния и ландшафт.
  \smallbreak
  \textbf{Оккультизм:} позволяет распознавать оккультную атрибутику, слова и концепции, народные традиции, а также опознавать древние тома с магией и оккультные коды. Помогает вспомнить тайные мистические знания, почерпнутые из старинных книг, учений или собственного опыта.
  \smallbreak
  \textbf{Оперировать Тяжелой Машинерией:} требуется для управления поездами, паровыми двигателями, бульдозерами или другими большими машинами.
  \smallbreak
  \textbf{Убеждение:} умение внушить человеку конкретную идею, концепцию или верование путем взвешенного аргументирования, дискутирования или обсуждения. Убеждение можно использовать не говоря правду. Успешное применение Убеждения требует времени, как минимум полчаса. Если вы хотите убедить кого-то на скорую руку - используйте Увещевание.
  \smallbreak
  \textbf{Пилотирование:} выберите специализацию, такую как Катер, Самолет или Дирижабль. Каждая специализация требует отдельной оплаты очками навыков. Навык отвечает за возможность безопасно управлять определенным видом транспорта.
  \smallbreak
  \end{multicols}
  \includegraphics[width=\linewidth]{img/bottom.png}
\end{fullcocpaperbox}

\begin{fullcocpaperbox}{}{}
  \raisebox{-10mm}{\includegraphics[width=\linewidth]{img/top.png}}
  \cocpaperboxtitle{Описание навыков}
  \begin{multicols}{3}
  \textbf{Психоанализ:} описывает владение разными видами психоэмоциональной терапии. Психоанализ может вернуть утраченные очки Рассудка: раз в игровой месяц, чтобы понять, насколько терапия помогла, подвергаемый ей персонаж должен кинуть 1D100 по значению навыка Психоанализа врача. Если бросок успешен, то пациент восстанавливает 1D3 очка Рассудка. Если неуспешен - ничего не происходит. Если бросок критически провален - тогда пациент теряет 1D6 очков Рассудка и лечение у текущего специалиста прекращается. В игре невозможно избавиться от неопределенного безумия (этот момент не описан в правилах Рассудка в стартере, к сожалению) - лечение от него требует пребывания 1D6 месяцев в психоневрологическом учреждении, на протяжении этого пребывания психотерапия будет лишь частью общего курса лечения. Успешное применение этого навыка так же может позволить персонажу на короткое время справиться с приступом фобии или мании или не попасть под власть иллюзий.
  \smallbreak
  \textbf{Психология:} эмпатия, которая есть у всех людей, позволяющая понять мотивы действий другого человека и его характер или осознать, когда он лжет. Хранитель может делать броски на Психологию в закрытую, сообщая правдивую или ложную информацию, как ее воспринимают персонажи.
  \smallbreak
  \textbf{Верховая Езда:} навык применяется к поездкам на взнузданных лошадях, ослах и мулах. Он гарантирует знание основ ухода за животным, снаряжением, и как управлять животным в галопе или на пересеченной местности. Если животное внезапно остановится или взбрыкнет, этот навык позволит наезднику остаться в седле.
  \smallbreak
  \textbf{Наука:} практические и теоретические познания в фундаментальных областях знания, полученные путем формального образования, хотя им так де могут владеть самообразованные ученые-любители. Масштаб знаний определяеися эпохой, в которой происходит игра. Очки навыков используются для получения специализации, например Астрономии, Биологии, Ботаники, Химии, Криптографии, Геологии, Фармакологии, Физики, Зоологии и т.п. Когда у персонажа нет конкретной специализации, они могут совершать броски по наиболее близкой специализации с увеличенной сложностью (или кубиком помехи), если того захочет Хранитель.
  \smallbreak
  \textbf{Ловкость Рук:} позволяет спрятать или замаскировать объект или объекты одеждой, лохмотьями и другими подходящими материалами. Также отвечает за аккуратность при управлении тонким оборудованием.
  \smallbreak
  \textbf{Наблюдательность:} навык поможет найти тайную дверь или отсек, заметить спрятавшегося воришку, увидеть неброскую улику, понять, что автомобиль был недавно перекрашен, понять, что рядом засада, заметить оттопырившийся карман и т.д. - один из наиболее важных навыков в арсенале исследователя. Когда исследователь ищет персонажа, который прячется, то сложность броска определяется показателем Скрытности оппонента.
  \smallbreak
  \textbf{Скрытность:} используется для избежания обнаружения, тихого передвижения и ухода от нежелательных лиц без привлечения их внимания.
  \smallbreak
  \textbf{Выживание:} знания, необходимые для выживания в экстремальных условиях, таких как пустыня, море, пустошь или полярная тундра. В этот навык входит умение охотиться, строить укрытия, понимание естественных угроз (какие ягоды есть можно, а какие - нет) и т.п. Очки навыков тратятся на определенные специализации, такие как Пустошь, Арктика, Пустыня, Море и т.д. Когда у персонажа нет конкретной специализации, они могут совершать броски по наиболее близкой специализации с увеличенной сложностью (или кубиком помехи), если того захочет Хранитель.
  \smallbreak
  \textbf{Плавание:} умение оставаться на поверхности и плыть в воде или другой жидкости. Совершайте бросок по этому навыку только в случае кризиса или опасности. Неудача на пушнутом броске на Плавание может привести к потере очков здоровья. Она так же может привести к тому, что персонажа утащит потоком или он вовсе утонет.
  \smallbreak
  \textbf{Бросок:} умение попадать по цели брошенным предметом. Предмет размером с ладонь может быть брошен на расстояние в метрах равное СИЛ персонажа, деленная на 5
  \smallbreak
  \textbf{Выслеживание:} умение следовать за человеком, транспортом или животным по земле или сквозь заросли. Факторы, влияющие на сложность: время, прошедшее с момента прохода того, кого выслеживают, дождь, тип поверхности и т.д.
  \end{multicols}
  \includegraphics[width=\linewidth]{img/bottom.png}
\end{fullcocpaperbox}

\end{document}